\section{Introduction}

Symmetry is among the most useful forms of structure in a multi-agent system.  Interchangeable agents, equivalent roles, repeated resources, and relabeled actions can collapse a large game into a much smaller quotient; symmetry can also guide equilibrium selection when agents must coordinate without a privileged labeling.  Recent work has formalized exact player--action relabelings in normal-form games, characterized the complexity of discovering them, and identified settings in which equilibria respecting the symmetries can be computed efficiently \cite{tewolde2025game}.  Exact equality, however, is brittle.  Nominally identical robots have different energy costs, estimated payoff tables contain noise, and game designers routinely introduce small role- or location-dependent perturbations.  In such cases an exact automorphism disappears even when the strategic interaction remains nearly unchanged.

A naive remedy is to measure how much payoffs change under relabeling.  Raw payoff error can be strategically misleading.  Adding an arbitrary function $c_i(a_{-i})$ to player $i$'s utility can make payoff tables arbitrarily far apart while leaving every unilateral incentive, best response, and Nash equilibrium unchanged.  Thus, a useful notion of approximate symmetry should compare the incentives that define strategic behavior rather than payoff levels themselves.

We study a finite game $\Gamma$ together with a candidate finite group $G$ of player and action relabelings.  Our objective is to turn $G$ into a \emph{certified abstraction}: project $\Gamma$ onto an exactly $G$-invariant surrogate $\widehat\Gamma$, solve or constrain the surrogate using its orbit structure, and translate the result back to the original game with an explicit regret certificate.  The central quantity is the minimum change in any unilateral-deviation incentive required to make $G$ exact.  We call this quantity the \emph{strategic symmetry defect} $\delta_G(\Gamma)$.

Our contributions are as follows.
\begin{itemize}
    \item We define approximate game symmetry through maximum pairwise difference (MPD), a seminorm that quotients out strategically irrelevant payoff components.  Unlike raw payoff distance, zero strategic defect means strategic equivalence to a symmetric game, not necessarily literal equality of payoff tables.
    \item We give a sparse linear program for the nearest exactly $G$-invariant surrogate.  More structurally, we show that the unrestricted projection is exactly a potential-fitting problem on a directed graph of symmetry orbits: $\delta_G$ equals the magnitude of a most-negative mean cycle.  The cycle is a concise local witness of why the proposed symmetry cannot be made exact at smaller error.
    \item We prove an additive equilibrium-transfer theorem: every $\epsilon$-Nash equilibrium of a surrogate at strategic distance $\delta$ is an $(\epsilon+\delta)$-Nash equilibrium of the original game.  It follows that every finite game admits a $G$-respecting $\delta_G$-equilibrium.  A cyclic family shows that the coefficient one is asymptotically tight.
    \item We analyze the Reynolds, or group-averaging, projection.  It is available in closed form from payoff-coordinate orbits and is always within a factor two of the optimal strategic projection.  Nested candidate groups therefore induce a certified trade-off between orbit compression and equilibrium fidelity.
    \item For two-player zero-sum games, we formulate an orbit-restricted linear program that directly minimizes saddle-point residual and obtain per-player and total exploitability certificates.  We also extend the guarantees to payoff tables estimated from bounded samples.
    \item Lightweight experiments on perturbed matrix games and heterogeneous role-assignment games verify the certificates, separate strategic from payoff-level asymmetry, and demonstrate the compression--fidelity and runtime trade-offs on a laptop.
\end{itemize}

\paragraph{Related work.}
Exact game symmetries have a long connection to symmetric equilibria, beginning with the existence arguments for symmetric games \cite{nash1951noncooperative}.  Tewolde et al.\ give a general framework for player and action relabelings, connect symmetry discovery to graph automorphism, and characterize the complexity of computing symmetry-respecting equilibria \cite{tewolde2025game}.  Recent complexity results further show that symmetric common-payoff games remain difficult despite their additional structure \cite{ghosh2026complexity}.  Our work takes a supplied group as input and asks a complementary robustness question: how much must unilateral incentives change before that group becomes exact?

The strategic metric is inspired by the maximum pairwise difference used to study near-potential games \cite{candogan2013near}.  There, MPD controls how closely learning dynamics resemble those in a potential game; here, it measures distance to the fixed subspace of a relabeling group and directly certifies equilibrium transfer.  Approximate permutation symmetry has also been studied in large-population stochastic and mean-field multi-agent learning \cite{yardim2025approximately}.  Our setting is finite normal-form games with arbitrary player--action relabelings, exact finite-sample certificates, and combinatorial projection algorithms.

Symmetry is widely exploited in multi-agent learning and zero-shot coordination.  Other-Play removes arbitrary conventions by optimizing performance under relabelings \cite{hu2020otherplay}, while subsequent work learns or exploits approximate symmetries when they are not known a priori \cite{muglich2022opl}.  These approaches optimize policies across transformed environments.  We instead certify a fixed game's strategic proximity to a proposed symmetry and bound the incentive loss from enforcing it.

Finally, our pipeline is related to state and game abstraction.  MDP homomorphisms preserve decision structure under aggregation \cite{ravindran2004mdp}, and extensive-form abstractions provide solution-quality bounds under imperfect merging \cite{kroer2014extensive}.  The present abstraction acts directly on players, actions, and payoff coordinates of a strategic game.  The maximum-mean-cycle characterization builds on classical difference-constraint and cycle-mean methods \cite{karp1978cycle}.
